\scanpage{218}
\deleted{Il lui correspondrait un point \deleted{$\unsure{e|\alpha}$} $\alpha$ de $\operatorname{Pic}(X/S)$, rat/$k(a)$, et un anneau local $\underline{\mathcal O}_\alpha$, \deleted{tel que} \deleted{\unsure{pour tout}} et un hom $\underline{\mathcal O}_a\unsure{-|\to}\underline{\mathcal O}_\alpha$, tel que pour \deleted{\unsure{tout}} \inserted{$\underline{\mathcal O}_a$-}algèbre \deleted{\unsure{locale}} \inserted{locale} $A'$, \deleted{\unsure{locale}}, [$\underline{\mathcal O}_{\unsure{x|a}}\to A'$ hom local], on a $\unsure{F|\mathscr F|\mathcal F}(A')_0\simeq\operatorname{Hom}$ des $\underline{\mathcal O}_a$-algèbres $(\underline{\mathcal O}_\alpha,A')$, où le $0$ en $F(A')$ désigne les \deleted{\unsure{hom}} éléments de $F(A')$ qui induisent l'élémt de degré $(-1,1)$ sur $F(A'/\unsure{\mathfrak r(A')|r(A')|\mathfrak m(A')})$. Comme il y en a exactement un, il faudrait donc que $\underline{\mathcal O}_{x}=\underline{\mathcal O}_\alpha$}

Considérons \deleted{l'ensemble $\mathbb Z\times\mathbb Z\amalg\mathbb Z\times U$} en effet le préschéma \deleted{suivant} $S_{\deleted{\unsure{?|\natural}}I}$ suivant [où $I$ est un ensemble d'indices arbitraire] : Comme ensemble, il est identique à $\deleted{\unsure{U\amalg I}}U\times\{I\}\amalg I$ ; pour tout $i\in I$, soit $S_i=U\times\{I\}\cup\{i\}$, donc on a $S_I=\bigcup_{i\in I}S_i$, et $S_i\cap S_j=U\times\{I\}$ \unsure{$\simeq$} en correspondance \unsure{canonique|biunivoque|biunivoquement} avec $U$. On Munit chaque $S_i$ de la structure de \deleted{\unsure{pré}} \inserted{schéma} affine obtenue par transport de structure à partir de $S$, et on recolle.
% archive sheet number: 217
